\documentclass[12pt]{article}
\usepackage[utf8]{inputenc}
\usepackage[T2A]{fontenc}
\usepackage[english,russian]{babel}
\usepackage{amsmath}
\usepackage{amssymb}
\usepackage[a4paper,margin=2cm]{geometry}
\usepackage{setspace}
\usepackage{hyperref}
\usepackage{mathtools}
\onehalfspacing
\sloppy

\title{Задача Коши. Глобальное и локальное условие Липшица. Теорема Пикара о существовании и единственности решения задачи Коши. Особые случаи задачи }

\begin{document}

\maketitle

\section{Примечание}
Данный файл основан не на наших лекциях, а на материалах из видео лекций
по дифференциальным уравнениям физического факультета МГУ. Ссылка на
плейлист: \url{https://www.youtube.com/playlist?list=PLcsjsqLLSfNB3Uu7EL56s-JclAjjuSVw_}.

\section{Скалярная Задача Коши}

Рассмотрим линейное однородное дифференциальное уравнение первого порядка:
\begin{equation}\label{eq:linear-first-order}
    y' + p(x)y = 0.
\end{equation}
Данное уравнение вместе с дополнительным условием
\begin{equation}\label{eq:initial-condition}
    y(x_0) = y_0,
\end{equation}
где $x_0$ и $y_0$ - некоторые числа, называется \textbf{задачей Коши}
для уравнения \eqref{eq:linear-first-order}.

Аналогичная запись:
\begin{equation}\label{eq:cauchy-problem}
    \begin{cases}
        y' + p(x)y = 0, \\
        y(x_0) = y_0.
    \end{cases}
\end{equation}
То есть задача Коши заключается в поиске одной из множества интегральных кривых уравнения \eqref{eq:linear-first-order},
проходящей через точку $(x_0, y_0)$.

\subsection{Решение задачи Коши}
Решение задачи Коши \eqref{eq:cauchy-problem} в таком случае будет иметь вид:
\begin{equation}\label{eq:cauchy-solution}
    y(x) = y_0 e^{\int_{x_0}^{x} p(t) dt}.
\end{equation}
Убедиться в этом можно, подставив \eqref{eq:cauchy-solution}
в \eqref{eq:linear-first-order} и проверив выполнение
начального условия \eqref{eq:initial-condition}.

Данное уравнение будет решением \eqref{eq:linear-first-order}
так как решение \eqref{eq:linear-first-order} в общем виде имеет вид:
\begin{equation}\label{eq:general-solution}
    y(x) = Ce^{\int p(x) dx},
\end{equation}
где $C$ - произвольная константа. Подставляя $y_0$ вместо $C$ получаем
равенство.

Для проверки условия \eqref{eq:initial-condition} подставляем $x = x_0$ в \eqref{eq:cauchy-solution}:
\begin{equation*}
    y(x_0) = y_0 e^{\int_{x_0}^{x_0} p(t) dt} = y_0 e^0 = y_0.
\end{equation*}
Получаем выполнение начального условия. Тем самым, решение задачи Коши
найдено.

\subsection{Особые случаи задачи Коши}
1. Рассмотрим случай когда $y(x_0) = 0$. То есть второе условие задачи Коши
является однородным уравнением. В таком случае решение задачи Коши
будет иметь вид:
\begin{equation*}
    y(x) = 0 \cdot e^{\int_{x_0}^{x} p(t) dt} = 0.
\end{equation*}
Таким образом, $y(x) = 0$ является решением задачи Коши
в случае однородного начального условия.

2. Рассмторим случай когда в точке $x_0$ функция $p(x)$ не определена.
В таком случае обычная постановка решения задачи Коши не имеет смысла,
так как через точку $x_0$ не может проходить интегральная кривая
уравнения \eqref{eq:linear-first-order}. Однако, если $p(x)$ имеет
устранимый разрыв в точке $x_0$, то решение
задачи Коши может быть продолжено в точку $x_0$. То есть задача может быть
переформулирована так: Найти решения вида $y = y(x)$ (или $x = x(y)$), обладающее свойством:
\begin{equation*}
    \lim_{x \to x_0} y(x) = y_0 (\text{или } \lim_{y \to y_0} x(y) = x_0).
\end{equation*}
то есть решение, примыкающее к точке $(x_0, y_0)$.

Например, рассмотрим
уравнение:
\begin{equation*}
    y' + \frac{1}{x}y = 0,
\end{equation*}
и задачу Коши с условием $y(0) = y_0$. В данном случае функция
$p(x) = \frac{1}{x}$ не определена в точке $x_0 = 0$. Однако,
решение задачи Коши может быть продолжено в точку $x_0 = 0$,
если $y_0 = 0$. В этом случае решение задачи Коши будет иметь вид:
\begin{equation*}
    y(x) = 0 \cdot e^{\int_{0}^{x} \frac{1}{t} dt} = 0.
\end{equation*}

\end{document}